凸集合保证两个可行方案之间可以安全移动；凸函数进一步保证，沿这条线段移动时，目标不会在中途鼓起一道隐藏山脊。二者合在一起，才构成凸优化的全局结构。

本单元依次回答四个问题：

\begin{enumerate}
\tightlist
\item
  弦在图像上方意味着什么从定义、Jensen 不等式和上境图建立凸性的几何意义。
\item
  切平面为何是全局下界说明可微凸函数的一阶条件，并把不可微点推广到次梯度。
\item
  曲率区分凸、严格凸与强凸用 Hessian 连接局部曲率与全局凸性，解释唯一性和算法速度来自何种额外结构。
\item
  怎样识别和构造凸函数汇总可靠的保凸运算，并通过 log-sum-exp、损失函数和透视函数展示如何组合复杂模型。
\end{enumerate}

读者不需要靠图形猜测所有函数。真正可复用的方法是：先从定义或一阶、二阶判据识别基本原子，再用已经证明的保凸规则构造复杂表达式。

\subsection{一个函数同时看三种判断}

令 \(f(x)=x^2\)。定义法：\(f(\theta x+(1-\theta)y)-[\theta f(x)+(1-\theta)f(y)] = \theta(1-\theta)(x-y)^2\ge0\)。一阶法：\(f(y)-f(x)-f'(x)(y-x) = (y-x)^2\ge0\)。二阶法：\(f''(x)=2>0\)。三种语言说的都是同一个事实：抛物线在弦下方。实际建模时，二阶判据对光滑函数最便捷；定义适合理论证明；一阶下界则直接给出切平面证书，通向次梯度和对偶。

\subsection{怎么读这一章}

这四篇是在用四种语言描述同一件事。弦给出定义，切平面给出一阶判断，曲率给出二阶判断，保凸运算则负责把简单例子拼成实际模型。第一次读时先弄清每种语言能回答什么。集合为何也必须凸，可回看上一单元 \hyperref[ch:072]{``凸集合：可行方案能否安全混合''}。

\subsection{本章篇目}

以下篇目按概念依赖排列。

\begin{itemize}
\tightlist
\item \hyperref[sec:09-Convex-Optimization/03-Convex-Functions/01]{``弦在图像上方意味着什么''}；
\item \hyperref[sec:09-Convex-Optimization/03-Convex-Functions/02]{``切平面为何是全局下界''}；
\item \hyperref[sec:09-Convex-Optimization/03-Convex-Functions/03]{``曲率区分凸、严格凸与强凸''}；
\item \hyperref[sec:09-Convex-Optimization/03-Convex-Functions/04]{``怎样识别和构造凸函数''}。
\end{itemize}
